\scanpage{235}
\mspage{108}
% archive sheet number: 234

\deleted{Notons d'ailleurs que pour prouver que \deleted{\unsure{toute}} $f\in\mathrm{TF}$ est dans $W_A$, et $=$ est
$W_A$-asphérique, on est ramené aussitôt au cas où son but est un objet $a$ de $A$,
\diagram{\begin{tikzcd}[ampersand replacement=\&]
 \deleted{\unsure{Y|F|X}}X\arrow[d] \& X_a\arrow[l]\arrow[d]\\
 Y \& a\arrow[l]
\end{tikzcd}}
donc l'objet final de la catégorie $A_{/a}$. On est ainsi ramené au cas d'une
catégorie $\deleted{\unsure{B=A_{/a}}}$ $B=\unsure{A_{/a}|A_a}$ avec objet final,
vu que l'hyp. $d)$ est stable par passage de $A$ à une localisée $A_{/a}$.

À prouver que dans le cas où $A$ a un objet final, donc est $W$-asphérique,
pour tout $X$ dans $A^\wedge$, $X\xrightarrow{f}e$ se factorise en
$X\hookrightarrow X'\to e$, avec $X\to X'$ mono et $X'\longrightarrow e$ dans $W_A$.
Je vois}

\unsure{Réellement|Brutalement|Décidément}, mon projet de démonstration se \unsure{modifie|modifia}, car bien que
$d)$ ait l'air très, très fort, comme condition,
\deleted{\illegible{1}} telle quelle je ne vois pas, \unsure{pratiquement|finalement},
comment en déduire $a)$. Mais en fait
\deleted{\unsure{alors}} je \unsure{peux|puis|vais} considérer que l'équivalence
de $b,c,d)$ est mise là pour mémoire, elle est établie dans loc cit,
donc je n'ai à me préoccuper que de l'implication $b\Longrightarrow a$, qui va résulter
